\documentclass{article}
\usepackage[T1]{fontenc}
\usepackage{lmodern}
\usepackage{graphicx}
\usepackage{amsmath,amssymb}
\usepackage[a4paper,margin=2cm]{geometry}
\usepackage[absolute,overlay]{textpos}
\setlength{\TPHorizModule}{1mm}
\setlength{\TPVertModule}{1mm}
\usepackage[indonesian]{babel}
\usepackage{hyperref}
\setlength{\emergencystretch}{2em}
% unit: Halaman 3

\begin{document}

% === Halaman 3 (python/native) ===

Contoh:

(i) 23 mod 5 = 3          karena 23 dibagi 5 = 4 + sisa 3

(ii) –41 mod 9 = 4        karena |-41| mod 9 = 5, dan 9 – 5 = 4

(iii) 17  2 (mod 3)        karena 3 | (17 – 2)   

(iv) –7  15 (mod 11)   karena 11 | (-7 – 15) 

(v) 23 dan 40 relatif prima sebab PBB(23, 4) = 1

(vi) 4–1 (mod 9)  7 (mod 9) karena 47  1 (mod 9)

(vii)23–1 (mod 10) = -3 (mod 10) karena 23(-3)  1 (mod 10)

Latihan: (a) Hitung -24 mod 11 = ?


    (b) 12–1 (mod 5)  ?

3

\end{document}
